% Part: first-order-logic
% Chapter: sequent-calculus
% Section: derivations

\documentclass[../../../include/open-logic-section]{subfiles}

\begin{document}

\iftag{FOL}
      {\olfileid{fol}{seq}{der}}
      {\olfileid{pl}{seq}{der}}

\olsection{\usetoken{P}{derivation}}

\begin{explain}
We hebben vastgelegd hoe een beginsequent eruitziet en de
afleidingsregels gegeven. !!^{derivation}s in de sequentenrekening
worden hieruit inductief opgebouwd: elke !!{derivation} bestaat ofwel
uitsluitend uit een beginsequent, ofwel uit één of twee
!!{derivation}s gevolgd door een afleidingsstap.
\end{explain}

\begin{defn}[$\Log{LK}$ !!{derivation}]
Een \emph{$\Log{LK}$-!!{derivation}} van een sequent~$S$ is een eindige
boom van sequenten die aan de volgende voorwaarden voldoet:
\begin{enumerate}
\item De bovenste sequenten van de boom zijn beginsequenten.
\item De onderste sequent van de boom is~$S$.
\item Elke sequent in de boom behalve $S$ is een premisse van een
  correcte toepassing van een afleidingsregel waarvan de conclusie
  direct onder die sequent in de boom staat.
\end{enumerate}
We noemen $S$ dan de \emph{eindsequent} van de !!{derivation} en zeggen
dat $S$ \emph{!!{derivable} is in $\Log{LK}$} (of
$\Log{LK}$-!!{derivable}).
\end{defn}

\begin{ex}
Elke beginsequent, bijvoorbeeld $!C \Sequent !C$, is !!a{derivation}.
Door bijvoorbeeld de regel $\LeftR{\Weakening}$ toe te passen,
verkrijgen we hieruit een nieuwe !!{derivation}:
\begin{prooftree}
\Axiom$ \Gamma \fCenter \Delta $
\RightLabel{\LeftR{\Weakening}}
\UnaryInf$ !A, \Gamma \fCenter \Delta$
\end{prooftree}
De regel is echter algemeen bedoeld: we mogen $!A$ in de regel door
elke !!{sentence} vervangen, dus bijvoorbeeld ook door~$!D$. Als de
premisse overeenkomt met onze beginsequent $!C \Sequent !C$, zijn zowel
$\Gamma$ als $\Delta$ gelijk aan~$!C$ en luidt de conclusie $!D, !C
\Sequent !C$. Het volgende is dus !!a{derivation}:
\begin{prooftree}
\Axiom$ !C \fCenter !C $
\RightLabel{\LeftR{\Weakening}}
\UnaryInf$ !D, !C \fCenter !C$
\end{prooftree}
Vervolgens kunnen we een andere regel toepassen, bijvoorbeeld
$\LeftR{\Exchange}$, waarmee twee !!{sentence}s aan de linkerkant worden
verwisseld. Het volgende is daarom eveneens een correcte
!!{derivation}:
\begin{prooftree}
\Axiom$ !C \fCenter !C $
\RightLabel{\LeftR{\Weakening}}
\UnaryInf$ !D, !C \fCenter !C$
\RightLabel{\LeftR{\Exchange}}
\UnaryInf$ !C, !D \fCenter !C$
\end{prooftree}
Bij deze toepassing van de regel, die de vorm
\begin{prooftree}
\Axiom$ \Gamma, !A, !B, \Pi \fCenter \Delta $
\RightLabel{\LeftR{\Exchange}}
\UnaryInf$ \Gamma, !B, !A, \Pi \fCenter \Delta,$
\end{prooftree}
heeft, waren $\Gamma$ en $\Pi$ beide leeg, was $\Delta$ gelijk aan
$!C$, en werden $!A$ en $!B$ respectievelijk ingevuld door $!D$
en~$!C$. Op vrijwel dezelfde manier zien we dat ook
\begin{prooftree}
\Axiom$ !D \fCenter !D $
\RightLabel{\LeftR{\Weakening}}
\UnaryInf$ !C, !D \fCenter !D$
\end{prooftree}
!!a{derivation} is. We kunnen deze twee !!{derivation}s nu met behulp
van $\RightR{\land}$ combineren. Die regel luidde
\begin{prooftree}
\Axiom$\Gamma \fCenter \Delta, !A$
\Axiom$ \Gamma \fCenter \Delta, !B$
\RightLabel{\RightR{\land}}
\BinaryInf$ \Gamma \fCenter \Delta, !A \land !B $
\end{prooftree}
In ons geval moeten de premissen overeenkomen met de laatste sequenten
van de !!{derivation}s die daarin uitkomen. Dit betekent dat $\Gamma$
gelijk is aan $!C, !D$, dat $\Delta$ leeg is, dat $!A$ gelijk is aan
$!C$ en dat $!B$ gelijk is aan $!D$. Wil de afleidingsstap correct
zijn, dan is de conclusie dus $!C, !D \Sequent !C \land !D$.
\begin{prooftree}
\Axiom$ !C \fCenter !C $
\RightLabel{\LeftR{\Weakening}}
\UnaryInf$ !D, !C \fCenter !C$
\RightLabel{\LeftR{\Exchange}}
\UnaryInf$ !C, !D \fCenter !C$
\Axiom$ !D \fCenter !D $
\RightLabel{\LeftR{\Weakening}}
\UnaryInf$ !C, !D \fCenter !D$
\RightLabel{\RightR{\land}}
\BinaryInf$ !C, !D \fCenter !C \land !D $
\end{prooftree}
We kunnen de premissen uiteraard ook omkeren; dan is $!A$ gelijk aan
$!D$ en $!B$ aan~$!C$.
\begin{prooftree}
\Axiom$ !D \fCenter !D $
\RightLabel{\LeftR{\Weakening}}
\UnaryInf$ !C, !D \fCenter !D$
\Axiom$ !C \fCenter !C $
\RightLabel{\LeftR{\Weakening}}
\UnaryInf$ !D, !C \fCenter !C$
\RightLabel{\LeftR{\Exchange}}
\UnaryInf$ !C, !D \fCenter !C$
\RightLabel{\RightR{\land}}
\BinaryInf$ !C, !D \fCenter !D \land !C $
\end{prooftree}
\end{ex}
\end{document}
